% =============================================================================
% Presentación Beamer P01 – Protección Ocular frente a Radiación Láser
% Técnicas de Medida – Ingeniería Aeronáutica
% Compilar con: latexmk -xelatex -interaction=nonstopmode presentacion_P01.tex
% =============================================================================
\documentclass[aspectratio=169,12pt]{beamer}

\usetheme{Madrid}
\usecolortheme{seahorse}
\usefonttheme{professionalfonts}

% ---- Encoding & Language ----
\usepackage{fontspec}
\usepackage[spanish]{babel}
\addto\captionsspanish{\renewcommand{\tablename}{Tabla}}

% ---- Math & Physics ----
\usepackage{amsmath,amssymb}
\usepackage{mathtools}
\usepackage{siunitx}
\sisetup{output-decimal-marker={,}}

% ---- TikZ ----
\usepackage{tikz}
\usetikzlibrary{babel}
\usetikzlibrary{arrows.meta,positioning,calc,shapes.geometric,decorations.markings}

% ---- Tables ----
\usepackage{booktabs}
\usepackage{colortbl}
\usepackage{tabularx}

% ---- Code ----
\usepackage{listings}
\lstset{
  basicstyle=\ttfamily\tiny,
  keywordstyle=\color{blue!80}\bfseries,
  commentstyle=\color{green!50!black},
  backgroundcolor=\color{gray!10},
  frame=single,
  breaklines=true
}

% ---- Custom Colors ----
\definecolor{aeroblue}{RGB}{0,51,102}
\definecolor{aerored}{RGB}{153,0,0}
\definecolor{aerogreen}{RGB}{0,102,51}
\definecolor{aeroyellow}{RGB}{204,153,0}
\definecolor{aerogreendark}{RGB}{0,80,40}
\definecolor{aerogreenlight}{RGB}{0,140,70}

% ---- GREEN THEME ----
\setbeamercolor{frametitle}{bg=aerogreen,fg=white}
\setbeamercolor{title}{bg=aerogreen,fg=white}
\setbeamercolor{structure}{fg=aerogreen}
\setbeamercolor{palette primary}{bg=aerogreen,fg=white}
\setbeamercolor{palette secondary}{bg=aerogreenlight,fg=white}
\setbeamercolor{palette tertiary}{bg=aerogreendark,fg=white}
\setbeamercolor{palette quaternary}{bg=aerogreendark,fg=white}
\setbeamercolor{section in toc}{fg=aerogreen}
\setbeamercolor{subsection in toc}{fg=aerogreendark}
\setbeamercolor{block title}{bg=aerogreen,fg=white}
\setbeamercolor{block body}{bg=aerogreen!10,fg=black}
\setbeamercolor{block title alerted}{bg=aerored,fg=white}
\setbeamercolor{block body alerted}{bg=aerored!10,fg=black}
\setbeamercolor{block title example}{bg=aerogreenlight,fg=white}
\setbeamercolor{block body example}{bg=aerogreenlight!10,fg=black}
\setbeamercolor{item}{fg=aerogreen}
\setbeamercolor{itemize item}{fg=aerogreen}
\setbeamercolor{enumerate item}{fg=aerogreen}

% ---- Title Info ----
\title[P01 -- Protección Láser]{P01: Protección Ocular frente a\\Radiación Láser}
\subtitle{Técnicas de Medida -- Ingeniería Aeronáutica}
\author{Laboratorio de Técnicas de Medida}
\date{}

% =============================================================================
\begin{document}

\begin{frame}
	\titlepage
\end{frame}

\begin{frame}{Índice}
	\tableofcontents
\end{frame}

% =============================================================================
\section{Introducción}
% =============================================================================

\begin{frame}{¿Por qué Seguridad Láser?}
	\begin{columns}
		\begin{column}{0.55\textwidth}
			\begin{itemize}
				\item Láseres: herramientas clave en ingeniería aeronáutica
				      \begin{itemize}
					      \item PIV (Velocimetría por Imagen de Partículas)
					      \item LDA (Anemometría Láser Doppler)
					      \item Metrología y ensayos no destructivos
				      \end{itemize}
				\item Radiación coherente, colimada y de alta irradiancia
				\item \textbf{Riesgo ocular irreversible} incluso con exposiciones breves
			\end{itemize}
		\end{column}
		\begin{column}{0.4\textwidth}
			\begin{tikzpicture}[scale=0.7]
				\draw[fill=aerogreen!20,draw=aerogreen,thick] (0,0) rectangle (1.5,0.6);
				\node at (0.75,0.3) {\tiny\bfseries LÁSER};
				\draw[aerored,ultra thick,->] (1.5,0.3) -- (4.5,0.3);
				\draw[aerored!30,thin] (1.5,0.3) -- (4.5,0.7);
				\draw[aerored!30,thin] (1.5,0.3) -- (4.5,-0.1);
				\draw[fill=gray!30] (4.5,-0.1) ellipse (0.25 and 0.45);
				\draw[aerogreen,very thick] (4.2,0) rectangle (4.8,0.5);
				\node[aerogreen,font=\tiny\bfseries] at (4.5,0.25) {OD};
			\end{tikzpicture}
		\end{column}
	\end{columns}
\end{frame}

\begin{frame}{Objetivos de la Práctica}
	\begin{enumerate}
		\item Identificar peligros de radiación láser en el laboratorio
		\item Aplicar clasificación de seguridad (IEC~60825 / ANSI~Z136.1)
		\item Calcular la Exposición Potencial ($H_0$) y la OD requerida
		\item Seleccionar EPO adecuado de un catálogo de \textbf{12 lentes}
		\item Resolver un caso de estudio asignado por equipo (4~equipos)
		\item Usar herramientas Python para automatizar y verificar
	\end{enumerate}
\end{frame}

% =============================================================================
\section{Fundamentos Teóricos}
% =============================================================================

\begin{frame}{Propiedades del Láser}
	\begin{columns}
		\begin{column}{0.5\textwidth}
			\begin{block}{4 Propiedades Clave}
				\begin{enumerate}
					\item \textbf{Monocromaticidad} \\ $\lambda$ única
					\item \textbf{Coherencia} \\ Ondas en fase
					\item \textbf{Direccionalidad} \\ Baja divergencia
					\item \textbf{Alta Irradiancia} \\ Gran energía en área reducida
				\end{enumerate}
			\end{block}
		\end{column}
		\begin{column}{0.45\textwidth}
			\begin{alertblock}{Consecuencia}
				$\rightarrow$ Concentración de $10^5$ en la retina para $\lambda = 400$--$1400$~nm
			\end{alertblock}
			\vspace{0.3cm}
			\begin{exampleblock}{Zona Crítica}
				Nuestros 3 láseres (514,5 / 532 / 635~nm) caen en la \textbf{zona visible de máximo riesgo retiniano}
			\end{exampleblock}
		\end{column}
	\end{columns}
\end{frame}

\begin{frame}{Clasificación de Seguridad Láser}
	\centering
	\small
	\begin{tabular}{clll}
		\toprule
		\textbf{Clase} & \textbf{Riesgo}                       & \textbf{Ejemplo}           & \textbf{Controles}        \\
		\midrule
		1              & Seguro normal                         & Sensores encapsulados      & Ninguno                   \\
		2              & Visible, reflejo palpebral            & Puntero                    & Señalización              \\
		\rowcolor{aeroyellow!30}
		3R             & Riesgo bajo directa                   & \textbf{Diodo alineación}  & Señalización + precaución \\
		3B             & Peligroso directa                     & LDA baja $P$               & EPO frecuente             \\
		\rowcolor{aerored!25}
		\textbf{4}     & \textbf{Directa + especular + difusa} & \textbf{PIV, LDA alta $P$} & \textbf{EPO obligatorio}  \\
		\bottomrule
	\end{tabular}

	\vspace{0.5cm}
	\begin{tikzpicture}
		\draw[->,very thick,aerogreen!60!black] (0,0) -- (10,0);
		\foreach \x/\lab in {0.5/1, 2/2, 4/3R, 6/3B, 8.5/4} {
				\fill[aerogreen!\x 0!red] (\x,0) circle (0.2);
				\node[above,font=\scriptsize\bfseries] at (\x,0.25) {\lab};
			}
		\node[font=\tiny,left] at (0,0) {Seguro};
		\node[font=\tiny,right] at (10,0) {Máximo peligro};
	\end{tikzpicture}
\end{frame}

\begin{frame}{Interacción Láser -- Ojo}
	\begin{tikzpicture}[scale=0.75]
		% Spectrum
		\fill[violet!40] (0,0) rectangle (1.5,0.8);
		\fill[yellow!40] (1.5,0) rectangle (4,0.8);
		\fill[red!20]    (4,0) rectangle (6.5,0.8);
		\node[font=\tiny,rotate=90] at (0.75,0.4) {UV};
		\node[font=\tiny,rotate=90] at (2.75,0.4) {Vis+IR-A};
		\node[font=\tiny,rotate=90] at (5.25,0.4) {IR-B/C};
		% Danger
		\draw[aerored,very thick] (0,1.2) -- (1.5,2.2) -- (2.75,3) -- (4,2) -- (6.5,1);
		\node[aerored,font=\tiny,above] at (2.75,3) {Máx. peligro retiniano};
		% Lasers
		\draw[blue!70,dashed] (2.4,0) -- (2.4,2.6);
		\node[blue!70,font=\tiny,left] at (2.4,2.6) {514,5};
		\draw[green!60!black,dashed] (2.6,0) -- (2.6,2.8);
		\node[green!60!black,font=\tiny,right] at (2.6,2.8) {532};
		\draw[red!70,dashed] (3.2,0) -- (3.2,2.5);
		\node[red!70,font=\tiny,right] at (3.2,2.5) {635};
		% Axis
		\draw[->] (0,0) -- (7,0) node[right,font=\tiny] {$\lambda$ (nm)};
		\node[font=\tiny] at (0,-0.2) {100};
		\node[font=\tiny] at (1.5,-0.2) {400};
		\node[font=\tiny] at (4,-0.2) {1400};
	\end{tikzpicture}

	\vspace{0.3cm}
	\begin{itemize}
		\item \textbf{UV (100--400~nm):} córnea/cristalino $\rightarrow$ fotoqueratitis
		\item \textbf{Visible + IR-A (400--1400~nm):} retina $\rightarrow$ \alert{daño irreversible}
		\item \textbf{IR-B/C ($>$1400~nm):} córnea $\rightarrow$ quemaduras
	\end{itemize}
\end{frame}

\begin{frame}{Parámetros Clave y Ecuaciones}
	\begin{columns}
		\begin{column}{0.5\textwidth}
			\begin{block}{Parámetros}
				\begin{itemize}
					\item $\lambda$ -- longitud de onda (nm)
					\item $P$ -- potencia CW (W)
					\item $E$ -- energía/pulso (J)
					\item $a$ -- diámetro haz (mm)
					\item MPE -- Exposición Máx. Permisible
				\end{itemize}
			\end{block}
		\end{column}
		\begin{column}{0.45\textwidth}
			\begin{block}{Ecuaciones}
				Área del haz:
				$A = \dfrac{\pi\,a^2}{4}$
				\vspace{0.2cm}

				Exposición potencial:
				$H_0 = \dfrac{P \text{ o } E}{A}$
				\vspace{0.2cm}

				\textbf{Densidad óptica:}
				\[\boxed{OD = \log_{10}\!\left(\frac{H_0}{H_{\mathrm{MPE}}}\right)}\]
			\end{block}
		\end{column}
	\end{columns}
\end{frame}

% =============================================================================
\section{Escenarios Láser}
% =============================================================================

\begin{frame}{Tres Escenarios de la Práctica}
	\centering
	\small
	\begin{tabular}{lccc}
		\toprule
		               & \textbf{Láser 1 (PIV)} & \textbf{Láser 2 (LDA)} & \textbf{Láser 3 (Alin.)} \\
		\midrule
		Tipo           & Nd:YAG $2\omega$       & Argón Ion CW           & Diodo CW                 \\
		$\lambda$ (nm) & 532                    & 514,5                  & 635                      \\
		Modo           & Pulsado                & CW                     & CW                       \\
		$E$ (mJ)       & 200                    & {--}                   & {--}                     \\
		$P$ (W)        & {--}                   & 1,5                    & 0,0045                   \\
		$\tau$ (ns)    & 8                      & {--}                   & {--}                     \\
		$f$ (Hz)       & 15                     & {--}                   & {--}                     \\
		$a$ (mm)       & 5                      & 1,2                    & 3                        \\
		\rowcolor{aerored!15}
		\textbf{Clase} & \textbf{4}             & \textbf{4}             & \textbf{3R}              \\
		\bottomrule
	\end{tabular}
\end{frame}

\begin{frame}{Catálogo de Gafas EPO (12 Lentes)}
	\centering
	\tiny
	\begin{tabular}{cllc}
		\toprule
		\textbf{ID} & \textbf{Lente}  & \textbf{Bandas de protección (OD @ $\lambda$)}       & \textbf{VLT} \\
		\midrule
		L-01        & Naranja         & OD~5+ @ 190--540~nm                                  & 40\%         \\
		L-02        & Verde           & OD~2+ @ 630--670~nm; OD~4+ @ 800--1100~nm            & 55\%         \\
		L-03        & Azul cobalto    & OD~7+ @ 190--400~nm; OD~6+ @ 532~nm; OD~5+ @ 1064~nm & 15\%         \\
		L-04        & Transparente    & OD~2+ @ 10\,600~nm (CO$_2$)                          & 90\%         \\
		L-05        & Multi-$\lambda$ & OD~7 @ 190--534~nm; OD~6 @ 800--820~nm               & 30\%         \\
		L-06        & Ámbar           & OD~3+ @ 190--532~nm                                  & 50\%         \\
		L-07        & Rojo oscuro     & OD~4+ @ 488--540~nm                                  & 35\%         \\
		L-08        & Gris            & OD~3+ @ 630--694~nm; OD~5+ @ 800--1100~nm            & 25\%         \\
		L-09        & Naranja claro   & OD~2+ @ 315--532~nm; OD~4+ @ 808--1100~nm            & 45\%         \\
		L-10        & Verde jade      & OD~6+ @ 190--540~nm                                  & 20\%         \\
		L-11        & Rosa            & OD~1+ @ 585--650~nm                                  & 60\%         \\
		L-12        & Multi-$\lambda$ & OD~4+ @ 190--534~nm; OD~5+ @ 630--700~nm             & 22\%         \\
		\bottomrule
	\end{tabular}
\end{frame}

\begin{frame}{Valores de MPE de Referencia}
	\begin{block}{MPE para los Escenarios Base}
		\begin{itemize}
			\item \textbf{L1} (532~nm, pulsado): $H_\mathrm{MPE} = 5{,}0\times10^{-7}$ J/cm$^2$
			\item \textbf{L2} (514,5~nm, CW 0,25~s): $H_\mathrm{MPE} = 2{,}5\times10^{-3}$ W/cm$^2$
			\item \textbf{L3} (635~nm, CW 0,25~s): $H_\mathrm{MPE} = 2{,}5\times10^{-3}$ W/cm$^2$
		\end{itemize}
	\end{block}
	\vspace{0.3cm}
	\begin{alertblock}{Exposición Prolongada ($>$10~s)}
		\begin{itemize}
			\item L2 (514,5~nm): $H_\mathrm{MPE} = 1{,}0\times10^{-3}$ W/cm$^2$
			\item L3 (635~nm): misma consideración aplica
		\end{itemize}
	\end{alertblock}
\end{frame}

% =============================================================================
\section{Cálculos de Protección}
% =============================================================================

\begin{frame}{Láser 1: PIV -- Nd:YAG 532~nm (Clase~4)}
	\begin{columns}
		\begin{column}{0.5\textwidth}
			\begin{block}{Cálculo}
				$a = 5~\text{mm} = 0{,}5~\text{cm}$
				\[A = \frac{\pi (0{,}5)^2}{4} = 0{,}1963~\text{cm}^2\]
				$E = 200~\text{mJ} = 0{,}2~\text{J}$
				\[H_0 = \frac{0{,}2}{0{,}1963} \approx 1{,}02~\text{J/cm}^2\]
				\[OD_\mathrm{req} = \log_{10}\!\left(\frac{1{,}02}{5\times10^{-7}}\right) \approx \mathbf{6{,}31}\]
			\end{block}
		\end{column}
		\begin{column}{0.45\textwidth}
			\begin{alertblock}{Evaluación (extracto)}
				\begin{itemize}
					\item L-01: OD~5 \textcolor{aerored}{\textbf{$\times$ Insuf.}}
					\item L-03: OD~6 \textcolor{aerored}{\textbf{$\times$ Insuf.}}
					\item L-05: OD~7 \textcolor{aerogreen}{\textbf{$\checkmark$ OK}}
					\item L-10: OD~6+ \textcolor{aeroyellow}{\textbf{? Marginal}}
				\end{itemize}
			\end{alertblock}
			\vspace{0.2cm}
			$\Rightarrow$ \textbf{L-05} (VLT 30\%)
		\end{column}
	\end{columns}
\end{frame}

\begin{frame}{Láser 2: LDA -- Ar-Ion 514,5~nm (Clase~4)}
	\begin{columns}
		\begin{column}{0.5\textwidth}
			\begin{block}{Cálculo}
				$a = 1{,}2~\text{mm} = 0{,}12~\text{cm}$
				\[A = \frac{\pi (0{,}12)^2}{4} = 0{,}01131~\text{cm}^2\]
				$P = 1{,}5~\text{W}$
				\[H_0 = \frac{1{,}5}{0{,}01131} \approx 132{,}6~\text{W/cm}^2\]
				\[OD_\mathrm{req} = \log_{10}\!\left(\frac{132{,}6}{2{,}5\times10^{-3}}\right) \approx \mathbf{4{,}72}\]
			\end{block}
		\end{column}
		\begin{column}{0.45\textwidth}
			\begin{exampleblock}{Evaluación (extracto)}
				\begin{itemize}
					\item L-01: OD~5 \textcolor{aerogreen}{\textbf{$\checkmark$ OK}} (VLT 40\%)
					\item L-05: OD~7 \textcolor{aerogreen}{\textbf{$\checkmark$ OK}} (VLT 30\%)
					\item L-10: OD~6 \textcolor{aerogreen}{\textbf{$\checkmark$ OK}} (VLT 20\%)
				\end{itemize}
			\end{exampleblock}
			\vspace{0.2cm}
			$\Rightarrow$ \textbf{L-01} (mejor VLT: 40\%)
		\end{column}
	\end{columns}
\end{frame}

\begin{frame}{Láser 3: Alineación -- Diodo 635~nm (Clase~3R)}
	\begin{columns}
		\begin{column}{0.5\textwidth}
			\begin{block}{Cálculo}
				$a = 3~\text{mm} = 0{,}3~\text{cm}$
				\[A = \frac{\pi (0{,}3)^2}{4} = 0{,}07069~\text{cm}^2\]
				$P = 4{,}5~\text{mW} = 0{,}0045~\text{W}$
				\[H_0 = \frac{0{,}0045}{0{,}07069} \approx 0{,}0637~\text{W/cm}^2\]
				\[OD_\mathrm{req} = \log_{10}\!\left(\frac{0{,}0637}{2{,}5\times10^{-3}}\right) \approx \mathbf{1{,}41}\]
			\end{block}
		\end{column}
		\begin{column}{0.45\textwidth}
			\begin{exampleblock}{Evaluación (extracto)}
				\begin{itemize}
					\item L-02: OD~2 \textcolor{aerogreen}{\textbf{$\checkmark$ OK}} (VLT 55\%)
					\item L-08: OD~3 \textcolor{aerogreen}{\textbf{$\checkmark$ OK}} (VLT 25\%)
					\item L-11: OD~1 \textcolor{aerored}{\textbf{$\times$ Insuf.}}
					\item L-12: OD~5 \textcolor{aerogreen}{\textbf{$\checkmark$ OK}} (VLT 22\%)
				\end{itemize}
			\end{exampleblock}
			\vspace{0.2cm}
			$\Rightarrow$ \textbf{L-02} (mejor VLT: 55\%)

			\vspace{0.2cm}
			\begin{block}{Nota}
				Clase~3R: protección recomendada pero no estrictamente obligatoria.
			\end{block}
		\end{column}
	\end{columns}
\end{frame}

\begin{frame}{Resumen de Resultados}
	\centering
	\begin{tabular}{lcccc}
		\toprule
		\textbf{Láser (Clase)}   & $\boldsymbol{\lambda}$         & $\boldsymbol{H_0}$
		                         & $\boldsymbol{OD_\mathrm{req}}$ & \textbf{Lente}                   \\
		\midrule
		\rowcolor{aerored!15}
		1 -- PIV (\textbf{4})    & 532~nm                         & 1,02 J/cm$^2$      & 6,31 & L-05 \\
		\rowcolor{aerored!10}
		2 -- LDA (\textbf{4})    & 514,5~nm                       & 132,6 W/cm$^2$     & 4,72 & L-01 \\
		\rowcolor{aeroyellow!15}
		3 -- Alin. (\textbf{3R}) & 635~nm                         & 0,064 W/cm$^2$     & 1,41 & L-02 \\
		\bottomrule
	\end{tabular}

	\vspace{0.5cm}
	\begin{tikzpicture}
		\node[draw=aerogreen,thick,rounded corners,fill=aerogreen!10,
			text width=10cm,align=center,font=\small] {
			\textbf{Regla:} OD\textsubscript{gafa} $\ge$ OD\textsubscript{req} a la $\lambda$ del láser.\\
			Si varias cumplen $\Rightarrow$ preferir la de mayor VLT.
		};
	\end{tikzpicture}
\end{frame}

\begin{frame}{Evaluación de Riesgo -- Escenarios Base}
	\centering
	\small
	\begin{tabular}{l c c c c l}
		\toprule
		\textbf{Escenario} & \textbf{Prob.}                                  & \textbf{Sev.}      & $\boldsymbol{R}$
		                   & \textbf{Nivel}                                  & \textbf{Controles}                    \\
		\midrule
		\rowcolor{aerored!15}
		PIV (Clase~4)      & Alto (4)                                        & Mayor (4)          & 16
		                   & \textcolor{aerored}{\textbf{Alto}}              & EPO + interlock                       \\
		\rowcolor{aeroyellow!15}
		LDA (Clase~4)      & Medio (3)                                       & Mayor (4)          & 12
		                   & \textcolor{aeroyellow!80!black}{\textbf{Medio}} & EPO + cortinas                        \\
		\rowcolor{aerogreen!15}
		Alin. (3R)         & MuyBajo (1)                                     & Menor (2)          & 2
		                   & \textcolor{aerogreen}{\textbf{Bajo}}            & Señalización                          \\
		\bottomrule
	\end{tabular}

	\vspace{0.4cm}
	\begin{block}{$R = \text{Probabilidad} \times \text{Severidad}$}
		Bajo ($<5$) \quad | \quad Medio (5--12) \quad | \quad Alto ($>12$)
	\end{block}
\end{frame}

% =============================================================================
\section{Casos de Estudio por Equipos}
% =============================================================================

\begin{frame}{Trabajo por Equipos: 4 Escenarios}
	\begin{block}{Instrucciones}
		Cada equipo recibe un láser diferente. Deben:
		\begin{enumerate}
			\item Calcular $A$, $H_0$ y $OD_{\mathrm{req}}$
			\item Evaluar las 12 lentes del catálogo
			\item Seleccionar la óptima (o justificar si ninguna cumple)
			\item \textbf{Evaluar el riesgo:} asignar $P$ (1--5), $S$ (1--5),
			      calcular $R = P \times S$ y ubicar en la matriz
			\item Resolver 3 problemas adicionales
			\item Verificar con Python / notebook
		\end{enumerate}
	\end{block}
\end{frame}

\begin{frame}{Equipo~1: Nd:YAG 1064~nm (Clase~4, Pulsado)}
	\begin{columns}
		\begin{column}{0.38\textwidth}
			\centering\scriptsize
			\begin{tabular}{ll}
				\toprule
				\textbf{Par.}      & \textbf{Valor}                    \\
				\midrule
				$\lambda$          & 1064~nm                           \\
				Modo               & Pulsado                           \\
				$E$                & 500~mJ                            \\
				$\tau$             & 10~ns                             \\
				$f$                & 10~Hz                             \\
				$a$                & 7~mm                              \\
				$H_{\mathrm{MPE}}$ & $5{,}0\!\times\!10^{-6}$~J/cm$^2$ \\
				\rowcolor{aerored!15}
				Clase              & \textbf{4}                        \\
				\bottomrule
			\end{tabular}
		\end{column}
		\begin{column}{0.28\textwidth}
			\begin{alertblock}{\scriptsize Riesgo}
				\scriptsize
				IR invisible, pulsado,\\
				alta energía.\\
				Evaluar $P$, $S$, $R$.
			\end{alertblock}
		\end{column}
		\begin{column}{0.3\textwidth}
			\begin{exampleblock}{\scriptsize Problemas}
				\scriptsize
				\begin{enumerate}
					\item[P1.1] $E$ a 100~mJ
					\item[P1.2] Refl. 60\%
					\item[P1.3] Doble $\lambda$
				\end{enumerate}
			\end{exampleblock}
		\end{column}
	\end{columns}
\end{frame}

\begin{frame}{Equipo~2: CO$_2$ CW a 10\,600~nm (Clase~4)}
	\begin{columns}
		\begin{column}{0.38\textwidth}
			\centering\scriptsize
			\begin{tabular}{ll}
				\toprule
				\textbf{Par.}      & \textbf{Valor} \\
				\midrule
				$\lambda$          & 10\,600~nm     \\
				Modo               & CW             \\
				$P$                & 50~W           \\
				$a$                & 8~mm           \\
				$H_{\mathrm{MPE}}$ & 0,1~W/cm$^2$   \\
				\rowcolor{aerored!15}
				Clase              & \textbf{4}     \\
				\bottomrule
			\end{tabular}
		\end{column}
		\begin{column}{0.28\textwidth}
			\begin{alertblock}{\scriptsize Riesgo}
				\scriptsize
				IR lejano, 50~W CW,\\
				daño corneal.\\
				Evaluar $P$, $S$, $R$.
			\end{alertblock}
		\end{column}
		\begin{column}{0.3\textwidth}
			\begin{exampleblock}{\scriptsize Problemas}
				\scriptsize
				\begin{enumerate}
					\item[P2.1] MPE corneal
					\item[P2.2] $P$ a 5~W
					\item[P2.3] Haz enfocado
				\end{enumerate}
			\end{exampleblock}
		\end{column}
	\end{columns}
\end{frame}

\begin{frame}{Equipo~3: Diodo 808~nm CW (Clase~3B)}
	\begin{columns}
		\begin{column}{0.38\textwidth}
			\centering\scriptsize
			\begin{tabular}{ll}
				\toprule
				\textbf{Par.}      & \textbf{Valor}                    \\
				\midrule
				$\lambda$          & 808~nm                            \\
				Modo               & CW                                \\
				$P$                & 300~mW                            \\
				$a$                & 2~mm                              \\
				$H_{\mathrm{MPE}}$ & $2{,}5\!\times\!10^{-3}$~W/cm$^2$ \\
				\rowcolor{aeroyellow!25}
				Clase              & \textbf{3B}                       \\
				\bottomrule
			\end{tabular}
		\end{column}
		\begin{column}{0.28\textwidth}
			\begin{alertblock}{\scriptsize Riesgo}
				\scriptsize
				IR cercano invisible,\\
				daño retiniano.\\
				Evaluar $P$, $S$, $R$.
			\end{alertblock}
		\end{column}
		\begin{column}{0.3\textwidth}
			\begin{exampleblock}{\scriptsize Problemas}
				\scriptsize
				\begin{enumerate}
					\item[P3.1] Bombeo 1064~nm
					\item[P3.2] VLT y ergonomía
					\item[P3.3] Exp. $>$10~s
				\end{enumerate}
			\end{exampleblock}
		\end{column}
	\end{columns}
\end{frame}

\begin{frame}{Equipo~4: He-Ne 632,8~nm (3R) + LDA Prolongada}
	\begin{columns}
		\begin{column}{0.45\textwidth}
			\centering\scriptsize
			\begin{tabular}{ll}
				\toprule
				\multicolumn{2}{c}{\textbf{Caso~A: He-Ne}}             \\
				\midrule
				$\lambda$          & 632,8~nm                          \\
				$P$                & 10~mW                             \\
				$a$                & 1,5~mm                            \\
				$H_{\mathrm{MPE}}$ & $2{,}5\!\times\!10^{-3}$~W/cm$^2$ \\
				\rowcolor{aeroyellow!20}
				Clase              & \textbf{3R}                       \\
				\midrule
				\multicolumn{2}{c}{\textbf{Caso~B: LDA prolongada}}    \\
				\midrule
				$H_{\mathrm{MPE}}$ & $1{,}0\!\times\!10^{-3}$~W/cm$^2$ \\
				\bottomrule
			\end{tabular}
		\end{column}
		\begin{column}{0.25\textwidth}
			\begin{alertblock}{\scriptsize Riesgo}
				\scriptsize
				A: visible, baja $P$.\\
				B: exposición larga.\\
				Evaluar ambos.
			\end{alertblock}
		\end{column}
		\begin{column}{0.27\textwidth}
			\begin{exampleblock}{\scriptsize Problemas}
				\scriptsize
				\begin{enumerate}
					\item[P4.1] 20~mW
					\item[P4.2] Dual $\lambda$
					\item[P4.3] Obturador
				\end{enumerate}
			\end{exampleblock}
		\end{column}
	\end{columns}
\end{frame}

% =============================================================================
\section{Controles de Seguridad}
% =============================================================================

\begin{frame}{Jerarquía de Controles}
	\begin{tikzpicture}[scale=0.8,every node/.style={font=\small}]
		% Pyramid layers
		\fill[aerogreen!80!black] (0,0) -- (5,0) -- (4.2,1) -- (0.8,1) -- cycle;
		\fill[aerogreen!60] (0.8,1) -- (4.2,1) -- (3.4,2) -- (1.6,2) -- cycle;
		\fill[aeroyellow!80] (1.6,2) -- (3.4,2) -- (2.8,3) -- (2.2,3) -- cycle;
		\fill[aerored!60] (2.2,3) -- (2.8,3) -- (2.5,3.8) -- cycle;
		% Labels
		\node[white,font=\small\bfseries] at (2.5,0.5) {Eliminación / Sustitución};
		\node[white,font=\small\bfseries] at (2.5,1.5) {Ingeniería};
		\node[font=\small\bfseries] at (2.5,2.5) {Admin.};
		\node[white,font=\small\bfseries] at (2.5,3.3) {EPP};
		% Right labels
		\node[right,font=\scriptsize,text width=5cm] at (5.5,0.5)
		{Usar láseres de clase inferior si es posible};
		\node[right,font=\scriptsize,text width=5cm] at (5.5,1.5)
		{Interlock, obturador, cabina, cortinas};
		\node[right,font=\scriptsize,text width=5cm] at (5.5,2.5)
		{LSO, SOPs, señalización, registro};
		\node[right,font=\scriptsize,text width=5cm] at (5.5,3.3)
		{Gafas EN~207 con OD y VLT adecuados};
		% Arrow
		\draw[->,very thick,gray] (5.8,4) -- (5.8,-0.3);
		\node[gray,font=\tiny,rotate=90] at (6.2,2) {Más efectivo};
	\end{tikzpicture}
\end{frame}

\begin{frame}{Matriz de Riesgos}
	\centering
	\begin{tikzpicture}[scale=0.65]
		\foreach \x in {0,...,4} {
				\foreach \y in {0,...,4} {
						\pgfmathsetmacro{\risk}{(\x+1)*(\y+1)}
						\pgfmathparse{\risk < 5 ? "aerogreen!40" : (\risk < 13 ? "aeroyellow!60" : "aerored!50")}
						\edef\fillcol{\pgfmathresult}
						\fill[\fillcol] (\x,\y) rectangle (\x+1,\y+1);
						\draw[gray] (\x,\y) rectangle (\x+1,\y+1);
						\node[font=\tiny] at (\x+0.5,\y+0.5)
						{\pgfmathprintnumber[fixed,precision=0]{\risk}};
					}
			}
		\foreach \x/\l in {0/MuyBajo,1/Bajo,2/Medio,3/Alto,4/MuyAlto} {
				\node[font=\tiny,rotate=45,anchor=west] at (\x+0.5,-0.1) {\l};
			}
		\foreach \y/\l in {0/Insignif.,1/Menor,2/Moder.,3/Mayor,4/Catastrófico} {
				\node[font=\tiny,anchor=east] at (-0.1,\y+0.5) {\l};
			}
		\node[font=\tiny\bfseries,rotate=90] at (-1.3,2.5) {Severidad};
		\node[font=\tiny\bfseries] at (2.5,-1) {Probabilidad};
		% Markers
		\fill[blue,opacity=0.8] (3.5,3.5) circle (0.2);
		\node[white,font=\tiny\bfseries] at (3.5,3.5) {PIV};
		\fill[blue!70,opacity=0.8] (2.5,3.5) circle (0.2);
		\node[white,font=\tiny\bfseries] at (2.5,3.5) {LDA};
		\fill[orange,opacity=0.8] (0.5,1.5) circle (0.2);
		\node[white,font=\tiny\bfseries] at (0.5,1.5) {Alin.};
	\end{tikzpicture}

	\vspace{0.2cm}
	\small PIV y LDA (\textbf{Clase~4}): zona roja. Alineación (\textbf{3R}): zona verde.
\end{frame}

% =============================================================================
\section{Herramientas de Software}
% =============================================================================

\begin{frame}[fragile]{Automatización con Python}
	\begin{columns}
		\begin{column}{0.55\textwidth}
			\begin{block}{Contenido del Repositorio}
				\texttt{Practicas/P01\_Proteccion\_Laser/}
				\begin{itemize}
					\item \texttt{src/P01\_Proteccion\_Laser.py}
					\item \texttt{notebooks/P01\_Proteccion\_Laser.ipynb}
					\item \texttt{data/epo\_lenses.csv}
				\end{itemize}
			\end{block}
		\end{column}
		\begin{column}{0.4\textwidth}
			\begin{exampleblock}{Funciones}
				\begin{itemize}
					\item \texttt{area\_cm2\_from...()}
					\item \texttt{h0\_exposure()}
					\item \texttt{od\_required()}
					\item \texttt{risk\_assessment()}
					\item \texttt{evaluate\_and\_plot()}
				\end{itemize}
			\end{exampleblock}
		\end{column}
	\end{columns}

	\vspace{0.3cm}
	\begin{lstlisting}[language=Python]
A = area_cm2_from_diameter_mm(5.0)  # 0.1963 cm^2
H0 = 0.2 / A                       # 1.02 J/cm^2
print(od_required(H0, 5e-7))       # 6.31
  \end{lstlisting}
\end{frame}

% =============================================================================
\section{Cuestionario y Formato}
% =============================================================================

\begin{frame}{Preguntas de Discusión}
	\begin{enumerate}
		\item ¿Qué significa «OD~5+ @ 190--540~nm»? ¿En cuánto reduce la intensidad?
		\item ¿Por qué un Clase~4 es peligroso con reflexiones difusas?
		\item ¿Qué 3 parámetros se necesitan para calcular $OD_\mathrm{req}$?
		\item Entre dos gafas con igual OD, una VLT~45\% y otra 20\%: ¿cuál elegir?
		\item ¿Gafas para 1064~nm protegen a 532~nm del mismo Nd:YAG?
		\item ¿Qué protocolo se sigue cuando ninguna lente del catálogo cumple?
	\end{enumerate}
\end{frame}

\begin{frame}{Formato del Informe (AIAA)}
	\begin{enumerate}
		\item \textbf{Resumen:} máx. 200 palabras
		\item \textbf{Introducción}
		\item \textbf{Fundamento Teórico:} propiedades, clasificación, MPE, OD
		\item \textbf{Metodología}
		\item \textbf{Resultados:} 3 escenarios base + caso del equipo + evaluación de riesgo
		\item \textbf{Discusión:} (incluir cuestionario)
		\item \textbf{Conclusiones:} 4--6 puntos
		\item \textbf{Referencias} (IEEE/AIAA)
		\item \textbf{Apéndice A:} gráficas del script Python
	\end{enumerate}
\end{frame}

% =============================================================================
\section{Conclusiones}
% =============================================================================

\begin{frame}{Conclusiones}
	\begin{enumerate}
		\item Catálogo de 12 lentes cubre escenarios de UV a IR lejano,
		      incluyendo casos donde \textbf{ninguna lente cumple}
		\item La selección de EPO requiere: cobertura a $\lambda$,
		      $OD \ge OD_{\mathrm{req}}$ y máxima VLT posible
		\item La duración de exposición modifica la MPE y puede cambiar
		      la lente requerida
		\item Jerarquía de controles: eliminación $>$ ingeniería $>$ admin. $>$ EPP
		\item Script Python automatiza $H_0$, OD y evaluación del catálogo
		\item Cada caso de estudio requiere evaluación de riesgo ($R = P \times S$)
		      con ubicación en la matriz 5$\times$5
	\end{enumerate}

	\vspace{0.5cm}
	\centering
	\begin{tikzpicture}
		\node[draw=aerogreen,very thick,rounded corners,fill=aerogreen!10,
			text width=8cm,align=center,font=\normalsize\bfseries] {
			Nunca trabaje con láseres Clase~3B/4\\sin protección ocular verificada.
		};
	\end{tikzpicture}
\end{frame}

\begin{frame}{}
	\centering
	\vspace{2cm}
	{\Huge\color{aerogreen}\bfseries ¿Preguntas?}\\[1cm]
	{\large P01 -- Protección Ocular frente a Radiación Láser}\\[0.5cm]
	{\small Técnicas de Medida -- Ingeniería Aeronáutica}
\end{frame}

\end{document}
